% !TeX root = ../main.tex
% ==============================================================================
% Physical Modeling Framework — Basis for Telemetry & Sensor Adaptation
% ==============================================================================

\section{Physical Modeling Framework}
\label{sec:physical_modeling}

To provide physically consistent telemetry for anomaly recognition, the kinematic simulation of the automated guided vehicle (AGV) is extended with an electromechanical powertrain model. The objective is not to replicate every microscopic effect of a physical robot, but to capture the dominant coupling between trajectory demands, mechanical loads, electrical responses, and thermal dissipation.

\subsection{Motor and Drivetrain Dynamics}

The traction drive is modeled as an equivalent permanent-magnet DC (PMDC) machine coupled to a reduction gearbox ($i_g$) and driven wheel ($r_w$).

The electrical armature circuit is governed by:
\begin{equation}
    V_m(t) = R_a(T_m) I_m(t) + L_a \frac{dI_m(t)}{dt} + K_e \omega_m(t)
    \label{eq:armature_base}
\end{equation}
where $V_m$ is the armature terminal voltage, $I_m$ the armature current, $R_a(T_m)$ the winding resistance, $K_e$ the back-EMF constant, and $\omega_m$ the rotor speed.

\paragraph{Quasi-static approximation ($L_a \approx 0$):}
Because the electrical time constant ($\tau_e = L_a / R_a \approx 40\,\mathrm{ms}$) is much shorter than the mechanical time constant ($\tau_m \approx 1\text{--}3\,\mathrm{s}$) and the state reporting rate, the inductive transient $L_a \frac{dI_m}{dt}$ is neglected. This yields the quasi-static relationship:
\begin{equation}
    V_m(t) = R_a(T_m) I_m(t) + K_e \omega_m(t)
\end{equation}

Electromagnetic torque follows:
\begin{equation}
    \tau_m(t) = K_t I_m(t)
    \label{eq:electromagnetic_torque}
\end{equation}
where $K_t$ is the torque constant. The mechanical torque balance at the motor shaft is:
\begin{equation}
    J_m \frac{d\omega_m(t)}{dt} + B_m \omega_m(t) = \tau_m(t) - \frac{\tau_L(t)}{\eta_g i_g} - \tau_{\mathrm{friction}}
    \label{eq:torque_balance}
\end{equation}
where $J_m$ is rotor inertia, $B_m$ viscous friction, $\eta_g$ gearbox efficiency, and $\tau_{\mathrm{friction}}$ unmodeled or degraded drivetrain friction. The wheel load torque $\tau_L(t)$ aggregates rolling resistance, road grade, vehicle acceleration, and drawbar forces:
\begin{equation}
    \tau_L(t) = r_w \left[ C_{rr} m_{\mathrm{total}} g \cos\theta + m_{\mathrm{total}} a(t) + m_{\mathrm{total}} g \sin\theta + F_{\mathrm{drawbar}} \right]
\end{equation}
where $m_{\mathrm{total}} = m_{\mathrm{chassis}} + m_{\mathrm{payload}}$ is the total vehicle mass.

\subsection{Kinematics and Wheel Slip}

The relationship between rotor speed and linear ground speed $v(t)$ accounts for transmission geometry and wheel slip:
\begin{equation}
    v(t) = \frac{\omega_m(t)}{i_g} r_w (1 - s(t))
\end{equation}
where $s(t) \in [0, 1)$ represents the slip ratio. Under nominal traction, $s(t) \approx 0$. When wheel slippage occurs, $s(t) \to 1$, which decouples wheel speed from chassis velocity ($v_{\mathrm{wheel}} \gg v_{\mathrm{ground}}$) while collapsing the ground reaction torque ($\tau_L \to 0$).

\subsection{Thermal Dissipation Models}

Winding resistance increases with copper temperature:
\begin{equation}
    R_a(T_m) = R_{a0} \left[ 1 + \alpha (T_m(t) - T_0) \right]
\end{equation}
where $R_{a0}$ is the baseline resistance at temperature $T_0$, and $\alpha$ is the copper temperature coefficient.

The thermal dynamics of the motor are modeled as a lumped first-order node accounting for Joule heating, core iron losses (Steinmetz formulation), and convective heat transfer to the ambient environment ($T_{\mathrm{amb}}$):
\begin{equation}
    C_{\mathrm{th}} \frac{dT_m(t)}{dt} = I_m(t)^2 R_a(T_m) + C_{fe} \omega_m(t)^{1.6} B_{\max}^2 - \frac{T_m(t) - T_{\mathrm{amb}}}{R_{\mathrm{th}}}
\end{equation}
Similarly, the battery pack temperature $T_{\mathrm{bat}}(t)$ accounts for internal resistive heating:
\begin{equation}
    C_{\mathrm{th,bat}} \frac{dT_{\mathrm{bat}}(t)}{dt} = I_{\mathrm{bat}}(t)^2 R_{\mathrm{int}} - \frac{T_{\mathrm{bat}}(t) - T_{\mathrm{amb}}}{R_{\mathrm{th,bat}}}
\end{equation}

\subsection{Battery Equivalent Circuit Model (ECM)}

Battery terminal voltage $V_{\mathrm{bat}}(t)$ is represented by an ohmic equivalent circuit:
\begin{equation}
    V_{\mathrm{bat}}(t) = V_{\mathrm{ocv}}(SOC) - I_{\mathrm{bat}}(t) R_{\mathrm{int}}(SOC, T_{\mathrm{bat}})
    \label{eq:battery_ecm}
\end{equation}
where open-circuit voltage $V_{\mathrm{ocv}}(SOC)$ scales with state of charge ($SOC$), and internal resistance $R_{\mathrm{int}}$ varies with $SOC$, temperature, and battery degradation. State of charge evolves via Coulomb counting based on total drawn power and charging currents:
\begin{equation}
    \frac{dSOC(t)}{dt} = -\frac{I_{\mathrm{bat}}(t)}{Q_{\mathrm{nom}}} \cdot 100
\end{equation}

\paragraph{Omission of the polarization branch ($V_{\mathrm{pol}}$):}
A complete 1-RC ECM introduces a capacitive polarization branch ($\frac{dV_{\mathrm{pol}}}{dt} = -\frac{V_{\mathrm{pol}}}{R_p C_p} + \frac{I_{\mathrm{bat}}}{C_p}$). This branch is deliberately omitted, and its transfer resistance $R_p$ is folded into $R_{\mathrm{int}}$ ($0.035 + 0.015 = 0.050\,\Omega$). 
\emph{Justification:} Detailed $R_p/C_p$ parameters cannot be calibrated without dedicated electrochemical pulse tests, and the polarization branch accounts for less than 20\% of total voltage sag. In a purely ohmic model, the load-step response is directly proportional to $I_{\mathrm{bat}} \cdot R_{\mathrm{int}}$, producing a clean, unambiguous signature for battery state-of-health (SOH) degradation without masking other anomalies.

\subsection{Physical Fault Injection and Transducer Emulation}

Rather than injecting synthetic noise into high-level features, faults perturb the physical equations directly:
\begin{itemize}
    \item \textbf{Mechanical Stall:} Imposes $\omega_m = 0$, collapsing back-EMF ($E_b = 0$) and causing the current to surge to the drive limit, accompanied by rapid thermal rise and battery voltage sag.
    \item \textbf{Wheel Slippage:} Forces $s(t) \to 1$, reducing load torque ($\tau_L \to 0$) while motor speed remains high.
    \item \textbf{Gearbox Degradation:} Injects an additive friction torque $\tau_{\mathrm{friction}}$, increasing steady-state current and operating temperature.
    \item \textbf{Stator Fault:} Injects harmonic current ripple and multiplies copper Joule heating.
    \item \textbf{Battery SOH Degradation:} Multiplies $R_{\mathrm{int}}$, producing pronounced voltage drops during acceleration steps.
\end{itemize}

Finally, physical states are conditioned through transducer models (Gaussian noise, ADC quantization, and thermal/gain drift) to generate realistic diagnostic telemetry for anomaly detection algorithms.
